\section{原理分析与硬件电路图}

\subsection{主功率电路}

主功率电路由前级三端口功率变换器和后级单相逆变器组成。前级变换器连接光伏输入、储能电池和直流母线，其控制目标随工作模式变化：当光伏功率充足时，优先跟踪光伏最大功率点并向负载供电，多余能量用于电池充电；当光伏功率不足时，电池端口补充能量以维持母线电压；当电池处于欠压、过温或禁止放电状态时，系统限制输出功率或进入保护状态。

后级逆变器采用单相全桥结构，将稳定的直流母线电压转换为 220 V/50 Hz 交流输出。逆变控制需要同时考虑输出电压幅值、频率稳定性、负载突变响应和直流母线扰动抑制。为降低前后级耦合带来的振荡风险，前级母线环带宽和后级逆变电压环带宽需要协调设计。

\begin{figure}[htbp]
  \centering
  \fbox{\parbox[c][3.8cm][c]{0.88\linewidth}{\centering
  此处插入主功率电路原理图：\\
  光伏端口、储能端口、三端口变换器、直流母线、单相逆变桥和交流滤波器
  }}
  \caption{主功率电路原理图占位}
  \label{fig:power_circuit}
\end{figure}

\subsection{功率器件选型与计算}

功率器件选型主要依据最大输入电压、母线电压、峰值电流、开关频率、导通损耗、开关损耗和散热条件确定。三端口前级器件需要重点关注高占空比工况下的电压应力和电流纹波；后级逆变桥器件需要重点关注母线电压裕量、交流负载冲击电流和死区引入的输出畸变。

功率开关管的基本损耗可以分为导通损耗和开关损耗：

\begin{equation}
  P_{\mathrm{cond}} = I_{\mathrm{rms}}^2 R_{\mathrm{ds(on)}}
\end{equation}

\begin{equation}
  P_{\mathrm{sw}} \approx \frac{1}{2} V_{\mathrm{ds}} I_{\mathrm{d}} (t_{\mathrm{on}} + t_{\mathrm{off}}) f_{\mathrm{s}}
\end{equation}

其中，$I_{\mathrm{rms}}$ 为开关管有效电流，$R_{\mathrm{ds(on)}}$ 为导通电阻，$V_{\mathrm{ds}}$ 和 $I_{\mathrm{d}}$ 分别为开关过程中的电压和电流，$t_{\mathrm{on}}$、$t_{\mathrm{off}}$ 为开通和关断时间，$f_{\mathrm{s}}$ 为开关频率。实际设计还需考虑驱动电阻、寄生电感、温度对导通电阻的影响以及 PCB 散热条件。

\todo{补充最终 MOSFET、二极管、电感、电容和采样器件型号，以及对应电压/电流裕量计算。}

\subsection{控制电路}

控制电路围绕三颗 MCU 展开。GD32G553 主控板需要提供多路 ADC 采样、PWM 输出、故障输入、通信接口和隔离/非隔离驱动接口。采样通道包括光伏电压、光伏电流、电池端口电压、电池端口电流、直流母线电压、逆变输出电压和逆变输出电流。PWM 输出包括三端口变换器驱动信号和单相逆变桥驱动信号。

GD32C113 BMS 板负责电池侧模拟前端、均衡控制和保护执行。若采用分立采样方案，需要保证单体电压采样精度和共模范围；若采用专用电池采样芯片，则 GD32C113 主要承担通信、状态机和保护逻辑。GD32VW553 通信板通过串口或 SPI 与主控交换运行数据和配置参数，并通过 Wi-Fi 将数据发送至上位机。

\begin{figure}[htbp]
  \centering
  \fbox{\parbox[c][3.5cm][c]{0.88\linewidth}{\centering
  此处插入控制系统框图：\\
  GD32G553 主控、GD32C113 BMS、GD32VW553 无线模块、采样、驱动与保护接口
  }}
  \caption{控制电路框图占位}
  \label{fig:control_circuit}
\end{figure}

\subsection{驱动电路}

驱动电路负责将 MCU PWM 信号转换为功率开关所需的栅极驱动信号。前级三端口变换器和后级逆变桥均需要合理设置驱动电压、栅极电阻、死区时间和欠压锁定阈值。对于高侧开关，需要考虑自举供电或隔离供电方式；对于逆变桥，需要防止上下桥臂直通，并在硬件层面预留快速关断路径。

驱动设计还需关注开关节点高 dv/dt 对采样和通信的干扰。PCB 布局中应尽量缩小栅极驱动回路和功率换流回路面积，驱动地与功率地连接应遵循单点或分区原则，并通过合理的 RC 吸收、TVS 或栅极负压关断策略降低误导通风险。

\subsection{采样电路}

采样电路是数字控制闭环的基础。电压采样一般采用电阻分压加 RC 滤波后送入 ADC，电流采样可采用采样电阻、霍尔传感器或隔离放大方案。母线电压和逆变输出电压采样需要满足高压隔离和安全间距要求，电池采样需要兼顾精度、温漂和保护响应速度。

ADC 采样时刻应与 PWM 同步，尽量避开开关瞬态噪声。对于电流环控制，采样延迟和滤波相位滞后会直接影响环路稳定性，需要在软件中配合数字滤波和限幅处理。关键保护量可以同时进入硬件比较器或外部保护芯片，以缩短严重故障下的关断时间。
